\section{Connected components of a simple normal crossings boundary}

The smooth members of a simple normal crossings boundary need not be
connected or irreducible, as in Bierstone--Milman\cite{bierstone-milman08-functoriality},
Definition 2.1. Multiplicities are therefore assigned to their genuine connected
components. The global component theorem of Demailly\cite{demailly-cadg},
\dcref{II.5.3}, supplies local finiteness; compactness is used only when selecting
a finite refinement.

\begin{definition}[The connected boundary components]
\label{hir:connected-boundary-components}
\uses{hir:snc}
\lean{Hironaka.HironakaConnectedBoundaryComponents.boundaryComponents,Hironaka.HironakaConnectedBoundaryComponents.sUnion_boundaryComponents,Hironaka.HironakaConnectedBoundaryComponents.boundaryComponents_eq_empty_of_isSNCDivisor_zero}
\leanok
For an ordered finite family \(E=[H_1,\ldots,H_s]\) of closed subsets of a
complex manifold, let \(\mathcal C(E)\) be the family of all nonempty connected
components of its members. Equivalently, its members are the sets
\(\operatorname{Conn}_{H_i}(x)\), where \(x\in H_i\); each set is included once
in the family. Then
\[
 \bigcup_{C\in\mathcal C(E)}C=\bigcup_i H_i.
\]
For an SNC family in ambient dimension zero, \(\mathcal C(E)\) is empty.
\end{definition}
\begin{proof}
Every connected component of \(H_i\) is contained in \(H_i\). Conversely,
if \(x\in H_i\), the component \(\operatorname{Conn}_{H_i}(x)\) contains
\(x\), and is one of the specified members. These two containments prove
the identity, including when the family is empty.
In ambient dimension zero, the simple normal crossings condition admits no
member through any point: its defining coordinate singleton would be a
singleton in the empty coordinate set. All members are therefore empty, and
the connected component family is empty as well.
\end{proof}

\begin{lemma}[The actual geometry of a boundary component]
\label{hir:connected-boundary-component-geometry}
\uses{hir:connected-boundary-components,rem:global-components}
\lean{Hironaka.HironakaConnectedBoundaryComponents.component_geometry}
\leanok
Let \(E\) be a simple normal crossings boundary in a complex manifold of
dimension \(m+1\). Every \(C\in\mathcal C(E)\) is a nonempty closed connected
complex submanifold of dimension \(m\). It is a component of some \(H_i\),
and at every \(x\in C\) there is an open neighbourhood \(U\) with
\(H_i\cap U=C\cap U\).
\end{lemma}
\begin{proof}
The normal crossings coordinates exhibit \(H_i\) locally as one coordinate
hyperplane; exchanging its defining coordinate with the last coordinate
exhibits dimension \(m\). Thus \(H_i\) is analytic and all its points are
regular. A connected component of \(H_i\) is closed: its closure is connected,
is contained in the closed set \(H_i\), and hence remains in the same maximal
connected component. The regular locus of \(H_i\) is all of \(H_i\), so its
global analytic components are precisely these already closed connected
components. At a regular point the connected component contains a sufficiently
small connected coordinate neighbourhood inside \(H_i\). This gives an
ambient open neighbourhood on which \(C\) and \(H_i\) agree, and transports
the hyperplane flattening to \(C\). The point defining the component supplies
its nonemptiness, and its connectedness is the defining maximality property.
\end{proof}

\begin{lemma}[Local finiteness and compact selection]
\label{hir:connected-boundary-components-locally-finite}
\uses{hir:connected-boundary-component-geometry,rem:global-components}
\lean{Hironaka.HironakaConnectedBoundaryComponents.exists_isOpen_finite_boundaryComponents,Hironaka.HironakaConnectedBoundaryComponents.locallyFinite_boundaryComponents,Hironaka.HironakaConnectedBoundaryComponents.finite_boundaryComponents_inter_compact}
\leanok
For a simple normal crossings boundary \(E\), the family \(\mathcal C(E)\)
is locally finite. For every compact set \(K\), only finitely many
\(C\in\mathcal C(E)\) meet \(K\).
\end{lemma}
\begin{proof}
The components of each \(H_i\) are its global analytic components, so the
global component theorem makes that family locally finite. At an arbitrary
point, intersect the finitely many open neighbourhoods supplied for the
original members \(H_i\). The resulting open neighbourhood meets only finitely
many of their components, and hence only finitely many members of
\(\mathcal C(E)\). To prove the compact assertion, choose such a neighbourhood
at each point of \(K\). A finite subcover exists. A component meeting \(K\)
meets one of these finitely many neighbourhoods, and therefore belongs to the
finite union of their finite meeting families.
\end{proof}

\begin{lemma}[Normal crossings survive finite component refinement]
\label{hir:snc-connected-component-refinement}
\uses{hir:connected-boundary-component-geometry,hir:snc}
\lean{Hironaka.HironakaConnectedBoundaryComponents.isSNCDivisor_component_refinement}
\leanok
Let \(F\) be a finite list of distinct members of \(\mathcal C(E)\), where
\(E\) is simple normal crossings. Then \(F\) is simple normal crossings.
\end{lemma}
\begin{proof}
Every member of \(F\) is closed by the component geometry lemma. Fix a point
\(x\) and normal crossings coordinates for \(E\). For each component in
\(F\) through \(x\), choose its original member \(H_i\) and an open
neighbourhood on which the two agree. Intersect these finitely many
neighbourhoods and restrict the coordinate chart to the intersection. The
selected components through \(x\) are now precisely the corresponding
coordinate hyperplanes. Two such components with the same original member
and containing \(x\) are the same connected component. Since the list has no
repetitions, their indices are equal. Distinct original members through
\(x\) have distinct defining coordinates by the original simple normal
crossings condition. The selected hyperplanes are therefore distinct, which
is the required condition for \(F\).
\end{proof}

\begin{theorem}[Finite refinement near a compact set]
\label{hir:compact-connected-boundary-refinement}
\uses{hir:connected-boundary-components-locally-finite,hir:snc-connected-component-refinement,hir:connected-boundary-components}
\lean{Hironaka.HironakaConnectedBoundaryComponents.exists_finite_refinement_on_compact,Hironaka.HironakaConnectedBoundaryComponents.exists_finite_refinement_of_support_subset_isCompact}
\leanok
For a simple normal crossings boundary \(E\) and a compact set \(K\), there
is a finite list \(F\), with no repetitions, whose members are exactly those
connected boundary components meeting \(K\). All its members are nonempty
closed connected smooth hypersurfaces, \(F\) is simple normal crossings, and
there is an open neighbourhood \(U\supset K\) such that
\[
 U\cap |F|=U\cap |E|,
 \qquad K\cap |F|=K\cap |E|.
\]
If \(|E|\subset K\), then \(|F|=|E|\).
\end{theorem}
\begin{proof}
List once each of the finitely many components meeting \(K\). The component
geometry and finite refinement lemmas give the stated smoothness and normal
crossings. The inclusion \(|F|\subset|E|\) follows from component containment.
The union of the omitted components is closed. Indeed, at a point outside
that union choose an open neighbourhood meeting only finitely many components,
and remove those finitely many closed omitted components from the neighbourhood.
The result is an open neighbourhood disjoint from the entire omitted union.
No omitted component meets \(K\). Its closed union therefore has an open
complement \(U\) containing \(K\). If \(x\in U\cap|E|\), choose an
original member containing \(x\). Its connected component through \(x\)
cannot be omitted, and thus appears in \(F\). This gives the support equality
on \(U\), and hence on \(K\). If the original support is contained in
\(K\), every nonempty component meets \(K\), so the same argument gives
the global support equality.
\end{proof}

\begin{lemma}[Divisibility depends on the local component]
\label{hir:boundary-component-germ-divisibility}
\uses{hir:connected-boundary-component-geometry,hir:divisor-exponent}
\lean{Hironaka.HironakaConnectedBoundaryComponents.isDivNear_connectedComponent_iff,Hironaka.HironakaConnectedBoundaryComponents.divExp_connectedComponent_eq}
\leanok
Let \(C\) be the connected component of \(H_i\) through a point of an SNC
boundary. At every \(x\in C\), a holomorphic germ ideal has the same local
divisibility exponents along \(C\) and along \(H_i\). In particular, their
maximal divisorial exponents at \(x\) are equal.
\end{lemma}
\begin{proof}
Choose an open neighbourhood \(U\) of \(x\) on which \(H_i\) and \(C\)
agree. For each integer \(r\), the defining local condition says that the
points of the carrier near \(x\) lie in the order-\(r\) cosupport. Intersecting
its neighbourhood with \(U\) makes the two conditions identical. The sets of
admissible powers therefore agree. Their maximal exponent, with the same
totalized convention when all powers are admissible, is equal.
\end{proof}

\begin{lemma}[Componentwise values glue locally constantly]
\label{hir:boundary-component-function-gluing}
\uses{hir:connected-boundary-component-geometry}
\lean{Hironaka.HironakaConnectedBoundaryComponents.isLocallyConstant_component_function}
\leanok
For any function \(a\) on \(\mathcal C(E)\), the function on \(H_i\) assigning
to \(x\) the value of \(a\) on the connected component of \(H_i\) through
\(x\) is locally constant.
\end{lemma}
\begin{proof}
Fix \(x\in H_i\) and its connected component \(C\). Choose an ambient open
neighbourhood \(U\) on which \(H_i\) and \(C\) agree. Every point of
\(H_i\cap U\) belongs to \(C\), and its connected component in \(H_i\)
is therefore the same set \(C\). The assigned value on this neighbourhood
is consequently constant. These relative neighbourhoods prove local
constancy at every point of \(H_i\).
\end{proof}

\begin{theorem}[Multiplicities on genuine boundary components]
\label{hir:connected-boundary-component-multiplicities}
\uses{hir:connected-boundary-component-geometry,hir:divisor-exponent-connected-carrier,hir:boundary-component-germ-divisibility,hir:boundary-component-function-gluing}
\lean{Hironaka.HironakaConnectedBoundaryComponents.exists_multiplicities_boundaryComponents,Hironaka.HironakaConnectedBoundaryComponents.exists_multiplicities_glued_boundaryComponents}
\leanok
Let \(D\) be an effective divisor and \(E\) a simple normal crossings
boundary. There is a function \(a:\mathcal C(E)\to\mathbb N\) such that,
for every \(C\in\mathcal C(E)\), \(x\in C\) and \(r\in\mathbb N\),
the ideal of \(D\) is locally divisible by \(\mathcal I_C^r\) at \(x\)
if and only if \(r\le a(C)\).
On every original member \(H_i\), the induced function assigning
\(a(\operatorname{Conn}_{H_i}(x))\) to \(x\) is locally constant, and
characterizes divisibility along \(H_i\) by the same inequality.
\end{theorem}
\begin{proof}
An invertible divisor ideal is coherent and has a nonzero stalk member at
every point. Its divisorial exponent along a smooth carrier is locally
constant. Each \(C\) is nonempty and connected, so this locally constant
function is constant on \(C\). Choose its value as \(a(C)\). The actual
maximal-exponent divisibility characterization gives the asserted equivalence
at every point. Choice of these values over the component family gives the
single function \(a\); there is no requirement that different connected
components of one original member receive the same value.
The component-function gluing lemma proves local constancy on the original
member. The local divisibility lemma transfers its characterization from
the connected component to that original member at each of its points.
\end{proof}

\begin{theorem}[Actual residual factorization after compact refinement]
\label{hir:compact-refined-boundary-residual}
\uses{hir:compact-connected-boundary-refinement,hir:divisor-finite-connected-component-residual}
\lean{Hironaka.HironakaConnectedBoundaryComponents.exists_residual_on_compact_SNC_components}
\leanok
Suppose the support of a simple normal crossings boundary \(E\) is contained
in a compact set. For every effective divisor \(D\), there is a finite
connected component refinement \(F\) with \(|F|=|E|\), a list of
coefficients \(b_C\), and an effective residual divisor \(R\) such that
\[
 \mathcal O(-D)=\left(\prod_{C\in F}\mathcal I_C^{b_C}\right)
                  \mathcal O(-R),\qquad
 \operatorname{Supp}(R)\subset\operatorname{Supp}(D),
 \qquad C\not\subset\operatorname{Supp}(R)\quad(C\in F).
\]
The coefficients in this factorization are extracted in order from successive
intermediate residual divisors.
\end{theorem}
\begin{proof}
Take the finite refinement of the preceding theorem. Each member is a
nonempty connected closed smooth hypersurface, so the finite connected
hypersurface residual theorem applies. Its induction successively removes
the maximal coefficient of the next member using the actual invertible colon
ideal. At each step the complete ideal factors as the chosen hypersurface
power times the new residual. Associating these equalities produces the
displayed product factorization. Support inclusion at each step preserves
noncontainment already obtained for earlier members, and gives the final
support inclusion. Empty refinements give the empty product and residual
\(D\), as required.
\end{proof}

\begin{remark}
On a noncompact manifold, a single disconnected smooth member can have
infinitely many connected components. The full component family is locally
finite; its finite selection near a compact set is not a global finite
refinement unless the original support is contained in that compact set.
An equality with a divisor sum additionally requires proving the residual
ideal is the unit ideal under the appropriate support hypothesis.
\end{remark}
